\section{Las opciones en las que participas}
\begin{itemize}
  \item Aprovechamiento académico
  \item Examen General de Egreso (EGEL)
  \item Créditos en estudios avanzados
  \item Artículo especializado para publicar en revista indizada
  \item Ensayo
  \item Memoria de experiencia laboral
  \item Reporte de aplicación de conocimientos
  \item Reporte de autoempleo profesional
  \item Reporte de residencia de investigación
  \item Tesina
  \item Tesis
\end{itemize}
\section{Lo que es igual en todas las opciones}
En todas eliges la opción en \emph{¿Cómo quieres titularte?}, capturas tus datos una vez, avanzas por etapas con la tarjeta \emph{Tu siguiente paso}, subes tus documentos en PDF y firmas tus formatos con un botón. Las pantallas de esta sección son de la opción \emph{Tesis}.
\captura{tesis_pasante_01_entrada}{«¿Cómo quieres titularte?»: abres «Con trabajo escrito» y eliges la opción.}
\captura{tesis_pasante_02_informativa}{Página de la opción, antes de iniciar: qué necesitas, cómo será el trámite y el manual para descargar.}
\captura{tesis_pasante_03_iniciar}{Aceptas el aviso de privacidad y pulsas «Iniciar mi trámite».}
\captura{tesis_pasante_07_captura_datos}{Formulario de la etapa: capturas tus datos una sola vez.}
\captura{tesis_pasante_13_subir}{Subir un documento: eliges el PDF y lo envías; queda «En validación».}
\captura{tesis_pasante_21_devuelto}{Un documento devuelto: arriba se muestra qué corregir y la observación de Titulación.}
\captura{tesis_pasante_22_resubir}{Corriges el documento devuelto: lees la observación, eliges el PDF corregido y lo vuelves a enviar.}
\captura{tesis_pasante_23_agenda1}{Eliges fecha y hora del examen previo y propones a tu sínodo; la solicitud llega al Coordinador de Docencia.}
\captura{tesis_pasante_40_estado}{«Estado de mi trámite»: en qué va y qué sigue, por etapas.}
\section{Lo propio de cada opción}
Requisitos y etapas de cada una, para compararlas. El detalle con pantallas está en el manual de cada opción.
\subsection{Aprovechamiento académico}
Reconoce el esfuerzo y la dedicación durante toda la trayectoria escolar. (Art. 14, REP). Sin trabajo escrito.

\textbf{Requisitos:}
\begin{itemize}
  \item Promedio general mayor o igual a 8.5.
  \item Haber cursado el plan de estudios sin interrupciones.
  \item Sin calificaciones reprobatorias ni anotaciones literales.
  \item Haber cubierto al menos el 50 \% de los créditos en la Universidad (si hubo revalidación o convalidación).
  \item Solicitar la opción por escrito dentro del primer año posterior a concluir la carrera.
  \item No haber sido sancionado por faltas a la responsabilidad universitaria.
\end{itemize}
\textbf{Etapas:}
\begin{enumerate}
  \item Tus datos
  \item Solicitud de la opción
  \item Documentación para la toma de protesta
  \item Fecha, hora y lugar de la toma de protesta
  \item Resultado de la evaluación
\end{enumerate}
\subsection{Examen General de Egreso (EGEL)}
Presentas el EGEL Plus del CENEVAL en sede UAEM y obtienes Desempeño Satisfactorio o Sobresaliente. (Art. 29, REP). Sin trabajo escrito.

\textbf{Requisitos:}
\begin{itemize}
  \item Tramitar la solicitud en el Departamento de Evaluación Profesional antes de presentar el examen.
  \item Presentar el EGEL en sede UAEM.
  \item Obtener testimonio de Desempeño Satisfactorio o Sobresaliente.
  \item Si el resultado es «Aún no satisfactorio» cuenta como evaluación aplazada; puedes pedir una segunda y última vez después de seis meses.
\end{itemize}
\textbf{Etapas:}
\begin{enumerate}
  \item Tus datos
  \item Solicitud para presentar el EGEL
  \item Registro y resultado del EGEL
  \item Documentación para la toma de protesta
  \item Fecha, hora y lugar de la toma de protesta
  \item Resultado de la evaluación
\end{enumerate}
\subsection{Créditos en estudios avanzados}
Acreditas módulos o unidades de un diplomado superior, especialidad o maestría afín a tu licenciatura. (Art. 22, REP). Sin trabajo escrito.

\textbf{Requisitos:}
\begin{itemize}
  \item Presentar la solicitud y obtener la autorización de la Subdirección Académica antes de inscribirte a los estudios.
  \item Que la Coordinación de Docencia dictamine la afinidad con tu licenciatura.
  \item Contar con carta de aceptación y con el plan de estudios oficial aprobado por los Consejos.
  \item Certificar el 100 \% de los créditos (diplomado superior o especialidad) o el 75 \% (maestría).
\end{itemize}
\textbf{Etapas:}
\begin{enumerate}
  \item Tus datos
  \item Solicitud de la opción
  \item Estudios avanzados cursados
  \item Documentación para la toma de protesta
  \item Fecha, hora y lugar de la toma de protesta
  \item Resultado de la evaluación
\end{enumerate}
\subsection{Artículo especializado para publicar en revista indizada}
Artículo original e inédito, enviado a una revista inscrita en índices validados. (Art. 16, REP). Con trabajo escrito.

\textbf{Requisitos:}
\begin{itemize}
  \item Artículo original e inédito, con autoría del pasante.
  \item Tema relacionado con el plan de estudios o con un área científica o técnica de la profesión.
  \item Evidencia de que la revista está en alguno de los índices del Art. 16 (e-Revist@s, Redalyc, Scielo, Scopus, CONACYT, Thomson-Reuters, JCR).
  \item Enviar el artículo hasta contar con el voto aprobatorio del asesor.
  \item Reporte de similitud iThenticate no mayor al 20 \%.
\end{itemize}
\textbf{Etapas:}
\begin{enumerate}
  \item Tus datos
  \item Protocolo y aprobación del tema
  \item Desarrollo del trabajo escrito
  \item Solicitud de revisión del trabajo escrito (examen previo)
  \item Solicitud de sustentación y voto aprobatorio
  \item Documentación para la evaluación profesional
  \item Solicitud del examen definitivo
  \item Resultado de la evaluación
\end{enumerate}
\subsection{Ensayo}
Exposición crítica y propositiva de un tema, con postura intelectual propia. (REP). Con trabajo escrito.

\textbf{Requisitos:}
\begin{itemize}
  \item Reporte de similitud iThenticate no mayor al 15 \%.
\end{itemize}
\textbf{Etapas:}
\begin{enumerate}
  \item Tus datos
  \item Protocolo y aprobación del tema
  \item Desarrollo del trabajo escrito
  \item Solicitud de revisión del trabajo escrito (examen previo)
  \item Solicitud de sustentación y voto aprobatorio
  \item Documentación para la evaluación profesional
  \item Solicitud del examen definitivo
  \item Resultado de la evaluación
\end{enumerate}
\subsection{Memoria de experiencia laboral}
Recopilación y resultados de una experiencia profesional en un ámbito reconocido. (Art. 33, REP). Con trabajo escrito.

\textbf{Requisitos:}
\begin{itemize}
  \item Constancia de la organización con al menos dos años continuos de servicios, o participación en un proyecto desde su inicio hasta su conclusión.
  \item El ingreso a los servicios o al proyecto debe ser posterior a la conclusión del plan de estudios.
  \item Reporte de similitud iThenticate no mayor al 20 \%.
\end{itemize}
\textbf{Etapas:}
\begin{enumerate}
  \item Tus datos
  \item Protocolo y aprobación del tema
  \item Desarrollo del trabajo escrito
  \item Solicitud de revisión del trabajo escrito (examen previo)
  \item Solicitud de sustentación y voto aprobatorio
  \item Documentación para la evaluación profesional
  \item Solicitud del examen definitivo
  \item Resultado de la evaluación
\end{enumerate}
\subsection{Reporte de aplicación de conocimientos}
Propuesta de solución a un problema determinado, derivada de lo aprendido en la carrera. (REP). Con trabajo escrito.

\textbf{Requisitos:}
\begin{itemize}
  \item Reporte de similitud iThenticate no mayor al 20 \%.
\end{itemize}
\textbf{Etapas:}
\begin{enumerate}
  \item Tus datos
  \item Protocolo y aprobación del tema
  \item Desarrollo del trabajo escrito
  \item Solicitud de revisión del trabajo escrito (examen previo)
  \item Solicitud de sustentación y voto aprobatorio
  \item Documentación para la evaluación profesional
  \item Solicitud del examen definitivo
  \item Resultado de la evaluación
\end{enumerate}
\subsection{Reporte de autoempleo profesional}
Informe sobre la constitución y desarrollo de una unidad económica formal. (Art. 52 y 53, REP). Con trabajo escrito.

\textbf{Requisitos:}
\begin{itemize}
  \item Unidad económica creada formal y legalmente (evidencia expedida por autoridad competente).
  \item Participación como propietario, socio, empresario o empleador durante al menos un año, concluido antes de iniciar el procedimiento.
  \item Individual o hasta tres pasantes del mismo programa.
  \item Reporte de similitud iThenticate no mayor al 20 \%.
\end{itemize}
\textbf{Etapas:}
\begin{enumerate}
  \item Tus datos
  \item Protocolo y aprobación del tema
  \item Desarrollo del trabajo escrito
  \item Solicitud de revisión del trabajo escrito (examen previo)
  \item Solicitud de sustentación y voto aprobatorio
  \item Documentación para la evaluación profesional
  \item Solicitud del examen definitivo
  \item Resultado de la evaluación
\end{enumerate}
\subsection{Reporte de residencia de investigación}
Informe de actividades de apoyo a la investigación bajo la dirección de un investigador activo en el SNI. (Art. 59, REP). Con trabajo escrito.

\textbf{Requisitos:}
\begin{itemize}
  \item Proyecto registrado en un centro o instituto de investigación de prestigio, con líneas afines a tu formación.
  \item Carta de estancia de al menos seis meses, expedida por la máxima autoridad del centro o instituto.
  \item Reporte de similitud iThenticate no mayor al 20 \%.
\end{itemize}
\textbf{Etapas:}
\begin{enumerate}
  \item Tus datos
  \item Protocolo y aprobación del tema
  \item Desarrollo del trabajo escrito
  \item Solicitud de revisión del trabajo escrito (examen previo)
  \item Solicitud de sustentación y voto aprobatorio
  \item Documentación para la evaluación profesional
  \item Solicitud del examen definitivo
  \item Resultado de la evaluación
\end{enumerate}
\subsection{Tesina}
Investigación documental de carácter monográfico, con aportación al estado del arte. (Art. 73, REP). Con trabajo escrito.

\textbf{Requisitos:}
\begin{itemize}
  \item Reporte de similitud iThenticate no mayor al 20 \%.
\end{itemize}
\textbf{Etapas:}
\begin{enumerate}
  \item Tus datos
  \item Protocolo y aprobación del tema
  \item Desarrollo del trabajo escrito
  \item Solicitud de revisión del trabajo escrito (examen previo)
  \item Solicitud de sustentación y voto aprobatorio
  \item Documentación para la evaluación profesional
  \item Solicitud del examen definitivo
  \item Resultado de la evaluación
\end{enumerate}
\subsection{Tesis}
Investigación documental, experimental, empírica o teórica; individual o de hasta tres pasantes. (Art. 80, REP). Con trabajo escrito.

\textbf{Requisitos:}
\begin{itemize}
  \item Extensión mínima de 60 cuartillas (individual) a interlineado máximo de 1.5.
  \item Reporte de similitud iThenticate no mayor al 20 \%.
\end{itemize}
\textbf{Etapas:}
\begin{enumerate}
  \item Tus datos
  \item Protocolo y aprobación del tema
  \item Desarrollo del trabajo escrito
  \item Solicitud de revisión del trabajo escrito (examen previo)
  \item Solicitud de sustentación y voto aprobatorio
  \item Documentación para la evaluación profesional
  \item Solicitud del examen definitivo
  \item Resultado de la evaluación
\end{enumerate}
